% s-165-p165-m003; original PDF 165, 166
\Aphorism{s-165-p165-m003}{252}
Man könnte in vielen ungleichartigen Lebenssphären dem Glücke nachgehen, wenn eins nicht wäre: das Netz von Beziehungsfäden, mit denen diese Sphären einander fassen und berühren. Dadurch ordnen sie sich einer gemeinschaftlichen Gesammtsphäre ein und müssen sich innerhalb dieser den Raum streitig machen. Traum und Wachen, Höhe und Niederung, Ekstase und Alltag, sie thun einander wehe, weil ihre Grenzlinie keine vollkommen scharfe, sondern eine empirisch getrübte, associativ überschreitbare ist, weil sie Einem und demselben Ich angehören, das sich nicht mathematisch theilen, sondern nur mechanisch zerstückeln kann und dabei, wie natürlich, Schmerz empfindet. Diese Krankheit legt sich mit den Jahren; am schlimmsten tritt sie in der Jugend auf, wo man noch gar keine getrennten Conti zu führen versteht und alles Mögliche, Religion, Liebe, Kunst, Wissenschaft, in eine hochgeschwollene Empfindungseinheit zusammendrängen möchte.
\par

